%==============================================================================
% REFERENSI (DAFTAR PUSTAKA)
% [Constraints]:
% - Gaya Sitasi: IEEE (nomor dalam kurung siku [1], [2], dst.)
% - Urutan: Alfabetis (abjad) berdasarkan nama pengarang
% - Kuota: Minimal 10 referensi (website maksimal 3, sisanya jurnal/buku/prosiding)
%==============================================================================
\chapter*{Referensi}
\addcontentsline{toc}{chapter}{Referensi}

\pet{Tuliskan berbagai referensi yang digunakan dalam laporan PKL terurut abjad berdasar nama pengarang dan beri nomor mulai dari [1], dengan gaya IEEE.}

\pet{Referensi yang digunakan \petb{minimal 10} (diperbolehkan mengambil referensi dari website, tetapi tidak boleh lebih dari 3).}

% Silakan buka komentar (uncomment) pada blok thebibliography di bawah ini untuk mengisi:
% Format IEEE:
% 1. Jurnal: Inisial. Nama, "Judul artikel," Nama Jurnal, vol. x, no. x, pp. xx-xx, Tahun.
% 2. Buku: Inisial. Nama, Judul Buku, Edisi ke-x. Kota Terbit: Penerbit, Tahun.
% 3. Prosiding: Inisial. Nama, "Judul makalah," dalam Nama Konferensi, Kota, Tahun, pp. xx-xx.
% 4. Website: Inisial. Nama / Organisasi, "Judul halaman web," URL, diakses Tanggal-Bulan-Tahun.

\begin{thebibliography}{99}
  \bibitem{ref1}
  A. B. Penulis, ``Penerapan Sistem Informasi Manajemen pada Rumah Sakit,'' \textit{Jurnal Rekayasa Sistem dan Teknologi Informasi}, vol. 5, no. 2, pp. 101--110, 2024.

  \bibitem{ref2}
  C. D. Author and E. F. Coauthor, \textit{Systems Analysis and Design Methods}, 7th ed. New York: McGraw-Hill, 2021.
\end{thebibliography}
